\section{Descent of morphisms of preschemes}
\label{VIII.5}

\begin{proposition}
\label{VIII.5.1}
Let $g\colon S'\to S$ be a morphism of preschemes.
\begin{enumerate}
\item[a)] Suppose that~$g$ is surjective, and that the homomorphism
$$
g^*\colon \mathcal{O}_{S}\to g_*(\mathcal{O}_{S'})
$$
is injective; then~$g$ is an epimorphism in the category
of preschemes, and even in the category of ringed
spaces.
\item[b)] Suppose that~$g$ is surjective and makes~$S$ a topological
quotient space of~$S'$. Let~$S''=S'\times_{S}S'$ and
let~$h\colon S''\to S$ be the structural morphism; consider the
canonical diagram of homomorphisms:
$$
\xymatrix@C=.5cm{\mathcal{O}_{S} \ar[r] &
g_*(\mathcal{O}_{S'})\ar@<2pt>[r]\ar@<-2pt>[r] &
h_*(\mathcal{O}_{S''})}\quoi,
$$
and suppose this diagram \emph{exact}. Then $g$ is an
effective epimorphism
in the category of preschemes (and also in the
category of ringed spaces), \ie the diagram
$$
\xymatrix@C=.5cm{S & S'\ar[l] & S''\ar@<2pt>[l]\ar@<-2pt>[l]}
$$
is exact.
\end{enumerate}
\end{proposition}

\subsubsection*{Proof} a) One must show that a morphism of ringed
spaces~$f\colon S\to Z$ is known when one
knows~$fg$. Now since~$g$ is surjective, one knows
the set-theoretic map~$f_{0}$ underlying~$f$; it remains to
determine the homomorphism of sheaves
of rings~$\mathcal{O}_{Z}\to \mathcal{O}_{S}$, or what amounts to the
same, the homomorphism
% original p. 211
of sheaves of rings
$$
u\colon f_{0}^{-1}(\mathcal{O}_{Z})\to \mathcal{O}_{S}
$$
defined by~$f$. One already knows the homomorphism
$$
(fg)_{0}^{-1}(\mathcal{O}_{Z})=
g_{0}^{-1}(f_{0}^{-1}(\mathcal{O}_{Z}))\to \mathcal{O}_{S'}
$$
defined by~$fg$, or what amounts to the same, one has at one's
disposal a homomorphism
$$
f_{0}^{-1}(\mathcal{O}_{Z})\to
g_{0*}(\mathcal{O}_{S'})=g_*(\mathcal{O}_{S'})\quoi.
$$
One observes at once that the latter is nothing but the composite
of~$g^*\colon \mathcal{O}_{S}\to g_*(\mathcal{O}_{S'})$ and
of~$u$, and since $g^*$ is injective, $u$ is known when one
knows~$g^*u$. [N.B. we have of course not used the fact
that~$g\colon S'\to S$ is a morphism of preschemes;
the statement would hold for an arbitrary morphism of ringed
spaces; the same remark applies to b), both in the
category of ringed spaces and in the category of
spaces ringed in local rings. Note also that if~$g$ is a
morphism of preschemes that is not necessarily surjective, but
such that~$g^*\colon \mathcal{O}_{S}\to g_*(\mathcal{O}_{S'})$
is injective, then for two morphisms~$f_{1},f_{2}$ from~$S$ into a
\emph{scheme}~$Z$ such that~$f_{1}g=f_{2}g$, one has $f_{1}=f_{2}$;
indeed, if~$I$ is the Ideal on~$S$ which defines the
subprescheme of~$S$ of coincidences of~$f_{1},f_{2}$
(inverse image of the diagonal subprescheme of~$Z\times Z$
under~$(f_{1},f_{2})$), one sees that~$I$ is contained
in~$\Ker(g^*)$].

b) One must show that for every ringed space~$Z$, the
following diagram of maps
$$
\xymatrix@C=.5cm{\Hom(S,Z)\ar[r] & \Hom(S',Z)\ar@<2pt>[r]\ar@<-2pt>[r] &
\Hom(S'',Z)}
$$
is exact, and that the same holds when~$Z$ is a space
ringed in local rings and one restricts oneself to homomorphisms
of spaces ringed in local rings. Since one already knows by a)
that the first map is injective, it remains to see that
if~$f'\colon S'\to Z$ is a homomorphism of ringed spaces such
that~$f'p_{1}=f'p_{2}$,
% original p. 212
then~$f'$ is of the form~$fg$, where~$f\colon S\to Z$ is a
homomorphism of ringed spaces. Since~$g$ is surjective, it is
then evident that if~$f'$ is a morphism of spaces ringed
in local rings, the same will hold for~$f$.

From the hypothesis on~$f'$ it follows that the underlying
set-theoretic map~$f_{0}'$ is constant on the fibers of
the map~$g_{0}$, hence, since the latter is surjective,
$f_{0}'$ factors uniquely as~$f_{0}'=f_{0}g_{0}$,
where~$f_{0}\colon S\to Z$ is a map, necessarily
continuous since~$g_{0}$ identifies~$S$ with a topological
quotient space of~$S'$. Now consider the homomorphism
$$
f_{0}^{-1}(\mathcal{O}_{Z})\to g_*(\mathcal{O}_{S'})
$$
deduced from the homomorphism~$(f_{0}g_{0})^{-1}(\mathcal{O}_{Z})\to
\mathcal{O}_{S'}$ corresponding
to~$f'$. The hypothesis~$f'p_{1}=f'p_{2}$ is then interpreted
as saying that the composites of the preceding homomorphism
with the two homomorphisms
$\xymatrix@C=.5cm{g_*(\mathcal{O}_{S'})\ar@<2pt>[r]\ar@<-2pt>[r] &
h_*(\mathcal{O}_{S''})}$ are the same, hence by
hypothesis b) it factors through a morphism
$$
f_{0}^{-1}(\mathcal{O}_{Z})\to \mathcal{O}_{S}\quoi.
$$
The latter defines a morphism of ringed spaces~$f\colon
S\to Z$, which is the desired morphism.

\begin{theorem}
\label{VIII.5.2}
Let $\cal{F}$ be the fibered category of arrows in the
category~$\Sch$ of preschemes
\textup{(VI~\Ref{VI.11}.a)}. Then every faithfully flat and
quasi-compact morphism~$g\colon S'\to S$ is an
$\cal{F}$-descent morphism (or again, as one says, a descent
morphism in~$\Sch$).
\end{theorem}

This therefore means the following: let~$S''=S'\times_{S}S'$, and for two
preschemes~$X,Y$ over~$S$, consider their inverse
images~$X',Y'$ over~$S'$ and their inverse images~$X'',Y''$
over~$S''$, whence a diagram of maps
$$
\xymatrix@C=.5cm{\Hom_{S}(X,Y)\ar[r] &
\Hom_{S'}(X',Y')\ar@<2pt>[r]\ar@<-2pt>[r] & \Hom_{S''}(X'',Y'')}\quoi;
$$
with these notations, \Ref{VIII.5.2} means that this diagram is
exact. Note that it
% original p. 213
is not true in general that~$g$ is an effective descent
morphism, \ie that for every prescheme~$X'$
over~$S'$, every descent datum on~$X'$ relative
to~$g\colon S'\to S$ is effective. The question of
effectiveness, often delicate, will be examined in
\No \Ref{VIII.7}.

We have seen \cite{VIII.D}, (taking into account that fibered products
exist in~$\Sch$) that statement~\Ref{VIII.5.2} is equivalent to the following:

\begin{corollary}
\label{VIII.5.3}
A faithfully flat and quasi-compact morphism of
preschemes is a universal effective epimorphism.
\end{corollary}

Since a faithfully flat and quasi-compact morphism remains such under
every extension of the base, one is reduced to proving that it is
an effective epimorphism. One then applies
criterion~\Ref{VIII.5.1} b), which gives the desired result, taking
into account~\Ref{VIII.4.3} and~\Ref{VIII.1.7}.

\begin{corollary}
\label{VIII.5.4}
Let $g\colon S'\to S$ be a faithfully flat and
quasi-compact morphism, $f\colon X\to Y$ an~$S$-morphism, $f'\colon X'\to Y'$
the~$S'$-morphism deduced from it by the change of base~$S'\to
S$. For~$f$ to be an isomorphism, it is necessary and sufficient that~$f'$
be one.
\end{corollary}

Indeed, if $f'$ is an isomorphism, it is also an isomorphism
for the natural descent structures on~$X',Y'$, and since the
functor~$X\mto X'$ from~$\Sch_{/S}$ into the category of
preschemes over~$S'$ endowed with a descent datum
relative to~$g$ is fully faithful by~\Ref{VIII.5.2},
the desired conclusion follows.

\begin{corollary}
\label{VIII.5.5}
Under the conditions of~\Ref{VIII.5.4}, for~$f$ to be a closed
immersion (\resp an open immersion, \resp a quasi-compact
immersion) it is necessary and sufficient that~$f'$ be one.
\end{corollary}

One may suppose as usual that~$Y=S$, and one has to prove
only the ``it is sufficient''. Note that the fact that~$X'/Y'$ is endowed
with a descent datum relative to~$g\colon Y'\to Y$, and
that the structural morphism~$f'\colon X'\to Y'$ is an immersion,
hence a monomorphism, implies that the two subobjects of~$Y''$
that are the inverse images
% original p. 214
of~$X'/Y'$ under the one and the other projection of~$S''$ into~$S'$
are the same. If~$f'$ is a closed immersion, it
follows by virtue of~\Ref{VIII.1.9} that there exists a
closed subprescheme $X_{1}$ of~$Y$ whose inverse image
under~$g\colon Y'\to Y$ is~$X'$. Hence by the uniqueness of the solution
of a descent problem relative to an
$\cal{F}$-descent morphism,
it
follows that~$X_{1}$ is $Y$-isomorphic to~$X$, hence $f\colon
X\to Y$ is a closed immersion. One proceeds in the same way for
an open immersion, using~\Ref{VIII.4.4}. If finally $f'$ is
a quasi-compact immersion, $f$ is quasi-compact by virtue
of~\Ref{VIII.3.3}, hence one can apply to the subset~$f(X)$
of~$Y$ the criterion~\Ref{VIII.4.5}, which proves that~$f(X)$ is
locally closed since its inverse image~$f'(X')$ in~$Y'$
is so. Replacing then~$Y$ by an open subset in
which~$f(X)$ is closed, one is reduced to the case where~$f'$
is a closed immersion, hence $f$ is one by virtue of what
precedes.

\begin{corollary}
\label{VIII.5.6}
Under the conditions of~\Ref{VIII.5.4}, for~$f$ to be affine, it
is necessary and sufficient that~$f'$ be so.
\end{corollary}

One proceeds as in~\Ref{VIII.5.5}, using~\Ref{VIII.2.1}
(One can also use Serre's cohomological criterion [EGA II
5.2], which proves~\Ref{VIII.5.6} without using any descent
technique).

\begin{corollary}
\label{VIII.5.7}
Under the conditions of~\Ref{VIII.5.4}, for~$f$ to be integral
(\resp finite, \resp finite and locally free) it is necessary and sufficient
that~$f'$ be so.
\end{corollary}

One has to prove only the ``it is sufficient'', and as usual
one may suppose $Y=S$, $Y$ affine, and $Y'$ affine. Since
the hypothesis implies that~$f'$ is affine, the same holds
for~$f$ by~\Ref{VIII.5.6}, hence~$X$ and consequently~$X'$ are
affine. Let~$A,A',B,B'=B\otimes_{A}A'$ be the rings
of~$Y,Y',X,X'$. One has $B=\varinjlim_i B_{i}$, where~$B_{i}$
runs through the sub-$A$-algebras of~$B$ that are of finite type
over~$A$, whence $B'=\varinjlim_i B_{i}'$, where
the~$B_{i}'$ are subalgebras of finite type of the
$A'$-algebra~$B'$. If $B'$ is integral
over~$A'$,
the~$B_{i}'$ are modules of finite type over~$A'$, hence, $A'$
being faithfully flat over~$A$, the~$B_{i}$ are modules of
finite type over~$A$, \ie $B$ is integral over~$A$. One sees in the same way
that if~$B'$ is finite over~$A'$, $B$ is so over~$A$. Same conclusion
for ``locally free of finite type'', \cf \Ref{VIII.1.11}.

\begin{corollary}
\label{VIII.5.8}
Under
% original p. 215
the conditions of~\Ref{VIII.5.4}, suppose~$f$ quasi-compact and
let~$\mathcal{L}$ be an invertible Module on~$X$, and~$\mathcal{L}'$
its inverse image on~$X'$. For~$\mathcal{L}$ to be ample
(\resp very ample) relative to~$f$, it is necessary and sufficient
that~$\cal{L}'$
be ample (\resp very ample) relative to~$f'$.
\end{corollary}

One has to prove only the ``it is sufficient''. The hypothesis
on~$\mathcal{L}$ implies in any case that~$f'$ is separated,
hence $f$ is separated by~\Ref{VIII.4.8}, and since~$f$ is
quasi-compact, and~$g\colon Y'\to Y$ is flat, the computation of direct
images by affine coverings shows that one has isomorphisms
$$
g^*(f_*(\mathcal{L}^{\otimes n}))\isomto
f_*'(\mathcal{L}'^{\otimes n})
$$
for every integer~$n$, hence one has an isomorphism
$$
g^*(\mathcal{S})\isomto \mathcal{S}'\quoi,
$$
where~$\mathcal{S}$ (\resp $\mathcal{S}'$) denotes the quasi-coherent
graded Algebra on~$Y$ (\resp on~$Y'$) that is the direct sum
of the~$f_*(\mathcal{L}^{\otimes n})$
(\resp of the~$f_*'(\mathcal{L}'^{\otimes n})$) for~$n\geq
0$. Note that for every~$n\geq 0$, the cokernel of the
canonical homomorphism~$f'^*(\mathcal{S}_{n}')\to \mathcal{L}'^{\otimes n}$
is the inverse image under~$X'\to X$ of the cokernel
of~$f^*(\mathcal{S}_{n})\to \mathcal{L}^{\otimes n}$, hence its
support~$Z_{n}'$ is the inverse image of the support~$Z_{n}$. If
$\cal{L}'$
is ample the intersection of the~$Z_{n}'$ is empty, hence, since~$X'\to X$
is surjective, the intersection of the~$Z_{n}$ is empty, \ie one has a
canonical morphism
$$
j\colon X\to \Proj(\mathcal{S})
$$
(EGA~II 3). Moreover, the analogous morphism
$$
j'\colon X'\to \Proj(\mathcal{S}')
$$
is none other than the one deduced from the preceding one by the
change of base~$Y'\to Y$ (\loccit). This being so, to say that
$\mathcal{L}'$ is ample relative to~$f'$ means that~$j'$ is
an immersion, necessarily quasi-compact moreover,
since~$f'$ is quasi-compact. Hence by virtue of~\Ref{VIII.5.5} $j$
is an immersion, \ie $\mathcal{L}$ is ample relative
to~$f$. --- One proceeds in an entirely analogous way in the case
of ``very
% original p. 216
ample'', restricting oneself above to~$n=1$, and replacing the
consideration of~$\Proj(\mathcal{S})$ by that of the projective
bundle~$\mathcal{P}(\mathcal{S}_{1})$ associated
to~$\mathcal{S}_{1}$.

Recall (EGA~II 5.1.1) that a quasi-compact morphism~$f$ is said to be
\emph{quasi-affine} if for every affine open set~$U$ in~$Y$,
$f^{-1}(U)$ is a prescheme isomorphic to an
open subscheme of an affine scheme. One shows (\loccit)
that it amounts to the same to say that~$\mathcal{O}_{X}$ is ample (or
also: very ample) relative to~$f$. Hence~\Ref{VIII.5.8}
implies:

\begin{corollary}
\label{VIII.5.9}
Under the conditions of~\Ref{VIII.5.4}, and supposing~$f$ quasi-compact,
for~$f$ to be quasi-affine, it is necessary and sufficient that~$f'$ be so.
\end{corollary}

\begin{remarks}
\label{VIII.5.10}
Hironaka's example of a non-projective variety shows that one
can have a proper morphism~$f\colon X\to Y$ of nonsingular
algebraic varieties (with~$Y$ projective), such that~$Y$
is the union of two open sets~$Y_{i}$ such
that~$X_{i}=X\times_{Y}Y_{i}$ is projective over~$Y_{i}$, without~$f$
being projective. Hence, setting~$Y'=Y_{1}\amalg Y_{2}$, $Y'$
is faithfully flat and quasi-compact (and even quasi-finite)
over~$Y$, $f'\colon X'\to Y'$ is projective, but~$f$ is not
projective. One must therefore be careful that in order to
apply~\Ref{VIII.5.8}, and to deduce from the fact that~$f'$ is
projective the same conclusion for~$f$, one must already have
on~$X'$ an invertible Module~$\mathcal{L}'$ ample for~$f'$,
\emph{endowed with a descent datum relative to}~$X'\to X$,
(which allows one to regard~$\mathcal{L}'$ as the inverse image
of an invertible Module~$\mathcal{L}$ on~$X$, which will then be ample
for~$f$ thanks to~\Ref{VIII.5.8}). When~$g\colon S'\to S$
is finite and locally free, see however~\Ref{VIII.7.7}.
\end{remarks}

\section[Application to finite and quasi-finite morphisms]{Application to finite and quasi-finite morphisms\protect\footnotemark}
\label{VIII.6}
\footnotetext{\Cf EGA IV 18.12 for generalizations to preschemes that are not necessarily locally noetherian}

We are going to prove the following two theorems:

\begin{theorem}
\label{VIII.6.1}
Let~$f\colon X\to Y$ be a \emph{proper morphism with finite fibers},
with~$Y$ locally noetherian. Then~$f$ is finite.
\end{theorem}

\begin{theorem}
\label{VIII.6.2}
Let~$f\colon X\to Y$ be a \emph{quasi-finite} and
\emph{separated} morphism, with~$Y$ locally noetherian. Then~$f$
is quasi-affine, and a fortiori quasi-projective.
\end{theorem}

\begin{remarks}
\label{VIII.6.3}
Theorem
% original p. 217
\Ref{VIII.6.1} is well known, and due to
Chevalley in the case of algebraic varieties; one will
also find a simple proof of it in [EGA III 4],
using the ``theorem on holomorphic functions''. The
proof given here does not use this latter
theorem, but on the other hand uses the theory of descent; we
give it as a bonus for the reader, since one gets it ``for nothing'' at
the same time as that of~\Ref{VIII.6.2}. Recall also ([EGA III
4] or \cite{VIII.1}) that the global form of Zariski's ``Main Theorem'',
deduced from the ``theorem on holomorphic functions'', asserts
that if~$f\colon X\to Y$ is quasi-finite and \emph{quasi-projective}, $Y$
being noetherian, then~$X$ is $Y$-isomorphic to an
open subprescheme of a
\emph{finite} $Y$-prescheme~$Z$. The conjunction of the ``Main Theorem'' and
of~\Ref{VIII.6.2} can therefore be stated as follows:
\end{remarks}

\begin{corollary}
\label{VIII.6.4}
Let~$f\colon X\to Y$ be a quasi-finite and separated morphism,
with~$Y$ noetherian. Then~$X$ is $Y$-isomorphic to an
open subprescheme of a finite $Y$-prescheme~$Z$.
\end{corollary}

Another interesting consequence of~\Ref{VIII.6.2} for the
theory of descent will be given with~\Ref{VIII.7.9}.

\subsubsection*{Proof of~\Ref{VIII.6.1} and~\Ref{VIII.6.2}}
We shall admit the following fact, whose proof is
easy\footnote{\Cf EGA IV 8}:

\begin{lemma}
\label{VIII.6.5}
Let~$X$ be a prescheme of finite type over~$Y$ locally
noetherian, and let~$y\in Y$. For there to exist an open
neighborhood~$U$ of~$y$ such that~$X|U$ is finite (\resp quasi-affine,
\resp\ ...) over~$U$, it is necessary and sufficient
that~$X\times_{Y}\Spec(\mathcal{O}_{y})$ be finite (\resp quasi-affine,
\resp\ ...) over~$\Spec(\mathcal{O}_{y})$.
\end{lemma}

Since on the other hand the property for~$f\colon X\to Y$
of being finite, \resp quasi-affine, is local on~$Y$, one is
reduced, in order to prove~\Ref{VIII.6.1} and~\Ref{VIII.6.2}, to the case
where~$Y$ is the spectrum of a local ring, and is a fortiori of
finite dimension. (N.B. one calls \emph{dimension of a
prescheme}~$Y$ the sup of the Krull dimensions of its local
rings). We proceed by induction on~$$
n=\dim(Y)\quoi,
$$
the assertion
% original p. 218
being trivial for~$n<0$. We may therefore suppose~$n\geq 0$,
and the assertion proved for the dimensions~$n'<n$. One may
again suppose that~$Y$ is the spectrum of a noetherian local
ring~$A$, of dimension~$n$. Note that the hypotheses
made in~\Ref{VIII.6.1} and~\Ref{VIII.6.2} are stable under
change of base (we have already used this in the reduction
at the beginning), so they will remain true after the change
of base~$\Spec(\widehat{A})\to \Spec(A)$. Since the latter is
faithfully flat and quasi-compact,
statements~\Ref{VIII.5.7} and~\Ref{VIII.5.9} reduce us to the
case where~$A$ is moreover complete. Using then the fact that every
noetherian local ring~$B$ over~$A$ that is quasi-finite over~$A$ is finite
over~$A$, and the fact that~$X$ is separated over~$Y$ and the fiber
of~$y$ consists of isolated points, one finds a
decomposition
$$
X=X'\amalg X''
$$
where~$X'$ is \emph{finite} over~$Y$, and where the fiber of~$X''$
at~$y$ is empty. If~$X$ is proper over~$Y$, so is~$X''$,
hence its image in~$Y$ is closed, and since it does not
contain~$y$, it is empty, hence~$X''=\emptyset$, hence~$X=X'$,
which shows that~$X$ is finite over~$Y$ and proves~\Ref{VIII.6.1}
(N.B. the induction hypothesis is useless here). If~$X$ is
quasi-finite over~$Y$, so is $X''$; now~$X''$ in fact lies over
the open set~$Y-\{ y\}$
of~$Y$, \emph{which is of dimension} $<n$. By virtue of the induction
hypothesis, $X''$ is quasi-affine
over~$Y-\{ y\}$,
hence also over~$Y$; the same obviously holds for~$X'$, hence
also for their sum~$X$, which proves~\Ref{VIII.6.2}.

\begin{remark}
\label{VIII.6.6}
Theorems~\Ref{VIII.6.1} and~\Ref{VIII.6.2} remain
valid if one no longer supposes~$Y$ locally noetherian,
provided one specifies that~$f$ is supposed to be of finite
presentation
(\cf \Ref{VIII.3.6}).
Indeed, one can still suppose~$Y$ affine, and then one verifies
without difficulty that the situation~$f\colon X\to Y$ is
deduced,
by a change of base~$Y\to Y_{0}$, from a situation~$f_{0}\colon
X_{0}\to Y_{0}$ satisfying the same hypotheses as~$f$,
with~$Y_{0}$ \emph{noetherian}. Hence by
result~\Ref{VIII.6.1} \resp \Ref{VIII.6.2}, $f_{0}$ is finite
\resp quasi-affine, hence so is~$f$. This kind of
reasoning is often useful for getting rid of
noetherian hypotheses,
% original p. 219
(which always end up being troublesome in the
applications).
\end{remark}
